\documentclass[11pt]{article}
\usepackage{amsmath,amssymb,amsthm}

\newtheorem{theorem}{Theorem}
\newtheorem{proposition}[theorem]{Proposition}

\begin{document}

\title{Completeness of the Divisor System}
\author{}
\date{}
\maketitle

\paragraph{System under consideration.}
For $a,r_2,r_3 \in \mathbb{Q}$ with $r_2r_3\neq 0$, define
\[
b:=\frac{a^2-r_2^2}{2r_2},\qquad
e:=\frac{a^2+r_2^2}{2r_2},\qquad
c:=\frac{a^2-r_3^2}{2r_3},\qquad
f:=\frac{a^2+r_3^2}{2r_3}.
\]
Then the two equations are
\[
a^2+b^2=e^2,\qquad a^2+c^2=f^2,
\]
equivalently
\[
a^2+\left(\frac{a^2-r_2^2}{2r_2}\right)^2
=
\left(\frac{a^2+r_2^2}{2r_2}\right)^2,
\]
\[
a^2+\left(\frac{a^2-r_3^2}{2r_3}\right)^2
=
\left(\frac{a^2+r_3^2}{2r_3}\right)^2.
\]

\begin{proposition}[Soundness]
For every $a,r_2,r_3\in\mathbb{Q}$ with $r_2r_3\neq 0$, the above formulas satisfy
\[
a^2+\left(\frac{a^2-r_2^2}{2r_2}\right)^2
=
\left(\frac{a^2+r_2^2}{2r_2}\right)^2,\qquad
a^2+\left(\frac{a^2-r_3^2}{2r_3}\right)^2
=
\left(\frac{a^2+r_3^2}{2r_3}\right)^2.
\]
\end{proposition}

\begin{proof}
For $i\in\{2,3\}$, clear denominators in
\[
a^2+\left(\frac{a^2-r_i^2}{2r_i}\right)^2
\stackrel{?}{=}
\left(\frac{a^2+r_i^2}{2r_i}\right)^2
\]
by multiplying by $4r_i^2$:
\[
4a^2r_i^2+(a^2-r_i^2)^2
\stackrel{?}{=}
(a^2+r_i^2)^2.
\]
Expanding both sides gives an identity, so both equations hold.
\end{proof}

\begin{theorem}[Completeness for integer solutions]
Let $a,b,e,c,f\in\mathbb{Z}$ satisfy
\[
a^2+b^2=e^2,\qquad a^2+c^2=f^2,
\]
and assume
\[
b<e,\qquad c<f.
\]
Then there exist $r_2,r_3\in\mathbb{Q}$ with $r_2\neq 0$, $r_3\neq 0$ such that
\[
b=\frac{a^2-r_2^2}{2r_2},\qquad
e=\frac{a^2+r_2^2}{2r_2},\qquad
c=\frac{a^2-r_3^2}{2r_3},\qquad
f=\frac{a^2+r_3^2}{2r_3}.
\]
In fact one can take
\[
r_2=e-b,\qquad r_3=f-c.
\]
\end{theorem}

\begin{proof}
Set $r_2:=e-b$ and $r_3:=f-c$. Since $b<e$ and $c<f$, we have $r_2,r_3\neq 0$.

From $a^2+b^2=e^2$:
\[
a^2=(e-b)(e+b)=r_2(e+b).
\]
Hence
\[
\frac{a^2-r_2^2}{2r_2}
=\frac{r_2(e+b)-r_2^2}{2r_2}
=\frac{e+b-r_2}{2}
=\frac{e+b-(e-b)}{2}
=b,
\]
and similarly
\[
\frac{a^2+r_2^2}{2r_2}
=\frac{e+b+r_2}{2}
=\frac{e+b+(e-b)}{2}
=e.
\]
The same argument with $(c,f,r_3)$ gives
\[
c=\frac{a^2-r_3^2}{2r_3},\qquad
f=\frac{a^2+r_3^2}{2r_3}.
\]
\end{proof}

\begin{theorem}[Completeness for rational solutions of the full three-equation system]
Let $a,b,c,e,f,g\in\mathbb{Q}$ satisfy
\[
a^2+b^2=e^2,\qquad
a^2+c^2=f^2,\qquad
b^2+c^2=g^2,
\]
and assume
\[
b<e,\qquad c<f.
\]
Then there exist $r_2,r_3\in\mathbb{Q}$ with $r_2\neq 0$, $r_3\neq 0$ such that
\[
b=\frac{a^2-r_2^2}{2r_2},\qquad
e=\frac{a^2+r_2^2}{2r_2},\qquad
c=\frac{a^2-r_3^2}{2r_3},\qquad
f=\frac{a^2+r_3^2}{2r_3},
\]
and
\[
\left(\frac{a^2-r_2^2}{2r_2}\right)^2
+
\left(\frac{a^2-r_3^2}{2r_3}\right)^2
=
g^2.
\]
In fact one can take
\[
r_2=e-b,\qquad r_3=f-c.
\]
\end{theorem}

\begin{proof}
Define $r_2:=e-b$ and $r_3:=f-c$. The strict inequalities give $r_2,r_3\neq 0$.

From the first two equations, exactly as above,
\[
b=\frac{a^2-r_2^2}{2r_2},\quad
e=\frac{a^2+r_2^2}{2r_2},\quad
c=\frac{a^2-r_3^2}{2r_3},\quad
f=\frac{a^2+r_3^2}{2r_3}.
\]
Now substitute the formulas for $b$ and $c$ into $b^2+c^2=g^2$:
\[
\left(\frac{a^2-r_2^2}{2r_2}\right)^2
+
\left(\frac{a^2-r_3^2}{2r_3}\right)^2
=
g^2.
\]
Hence every rational solution of the full three-equation system (under
$b<e$, $c<f$) arises from the divisor parameters $r_2,r_3$.
\end{proof}

\paragraph{Conclusion.}
Under the orientation assumptions $b<e$ and $c<f$ used in the Lean formalization
(\texttt{divisor\_system\_complete\_integer}), the parametrized system is complete:
it encompasses all integer solutions of
\[
a^2+b^2=e^2,\qquad a^2+c^2=f^2.
\]

\end{document}
